\section{Static estimates and Gaussian observation}\label{sec:static}
\providecommand{\sech}{\operatorname{sech}}

Our goal is a uniform heat-kernel bound whose logarithm, divided by $n$,
can be made small by choosing a long but fixed time.  We obtain it in four
steps: control the mean under an arbitrary change of law; transport that
control along a Gaussian observation with simultaneous posterior and cavity
gaps; bound entropy at the product endpoint; and convert the entropy bound
to a semigroup norm estimate.  The first step strengthens the stable-field
estimates of \cite{SKGap}.  The observation step uses
the stochastic-localization method of Eldan~\cite{Eldan2013}, its
Gaussian-channel formulation by El Alaoui and Montanari~\cite{ElAlaouiMontanari2021},
and the localization-annealing framework of Chen and Eldan~\cite{CE2025}.
The stopped, changed-law estimates and simultaneous cavity conclusions
required here are proved below. Throughout
this section constants may depend on the fixed $\beta<1$.  They never depend
on the probability measure being compared with equilibrium.

The cost of changing a strictly positive cube law $P$ to a probability
$Q$ will be measured by its square-root Dirichlet energy and relative
entropy.  Set
\[
 D_P(Q)=\mathcal E_P\bigl(\sqrt{Q/P}\bigr),\qquad
 H_P(Q)=\sum_x Q(x)\log\frac{Q(x)}{P(x)}.
\]
The energy lies in $[0,n]$; values at zeros of $Q$ are defined by
continuity.  To see how it controls the change in a single conditional
spin law, write $P_i,Q_i$ for the conditional laws of spin $i$ given
all other spins.  Expanding the conditional variance gives
\begin{equation}\label{eq:stat-affinity}
 D_P(Q)=\sum_i\E_{Q_{-i}}\left[
 1-\left(\sum_{s=\pm1}\sqrt{P_i(s)Q_i(s)}\right)^2\right].
\end{equation}
Conditionals on a $Q_{-i}$-null configuration can be assigned arbitrarily.
For two probabilities $p,q$, Cauchy--Schwarz applied to
$(\sqrt p-\sqrt q)(\sqrt p+\sqrt q)$ gives
$\TV(p,q)^2\le1-(\sum\sqrt{pq})^2$.  Hence, if
$\bar m_i=\E_Q[x_i\mid x_{-i}]$ and
$m_i^P=\E_P[x_i\mid x_{-i}]$, then
\begin{equation}\label{eq:stat-conditional-means}
 \E_Q\sum_i(\bar m_i-m_i^P)^2\le4D_P(Q).
\end{equation}

\subsection{A stable field controls means under an arbitrary law}

We begin with the deterministic hypotheses under which the mean estimate
will hold.  These hypotheses come from \cite{SKGap}; the arbitrary-law
conclusion will be proved here.  For a symmetric interaction matrix $M$
with zero diagonal,
a parameter $0<j<1$, and a field $h$, put
\[
 m(y)=\tanh y,\qquad b(y)=\frac1n\sum_i\sech^2 y_i,
 \qquad V(y)=\diag(\sech^2 y_i),
\]
and, for a nonnegative diagonal matrix $A$, write
\begin{align}
 F_{M,h}(y)&=y-h+\{j b(y)I-M\}m(y),\label{eq:stat-field-map}\\
 \mathsf H_M(y,A)
 &=I+A^{1/2}\left(jb(y)I-M-\frac{2j}{n}
                 m(y)m(y)^{\mathsf T}\right)A^{1/2}.
 \label{eq:stat-field-hessian}
\end{align}
The stability hypothesis is a uniform implication over a small-residual
region and nearby diagonals.  More precisely, the field has fixed positive
margins $\rho_0,\epsilon_A,\kappa$ and $a_+>1$, with $j a_+^2<1$, such that
\begin{equation}\label{eq:stat-stability}
 \left.
 \begin{gathered}
 \norm{F_{M,h}(y)}\le\rho_0\sqrt n,\quad 0\le A\le a_+I,\\
 \norm{A-V(y)}_{\HS}\le\epsilon_A\sqrt n
 \end{gathered}\right\}
 \quad\Longrightarrow\quad \mathsf H_M(y,A)\succeq\kappa I.
\end{equation}
We also require the finite-word estimates of \cite[Lemma~3.1]{SKGap}.
Fix a length bound $L$, a diagonal norm bound $M_0$, and an inverse margin
$c_0>0$. A word $W$ is an ordered product of at most $L$ factors, each
either a real diagonal matrix of norm at most $M_0$, a copy of $M$, or
the following inverse, which may occur at most once:
\[
 K_A=(I-AM+j\chi_AA)^{-1},\qquad \chi_A=n^{-1}\Tr A,
 \qquad 0\le A\le a_+I.
\]
Words containing $K_A$ are considered only when
\[
 I+j\chi_AA-A^{1/2}MA^{1/2}\succeq c_0I.
\]
To define the diagonal prediction $\mathsf P(W)$ for a word without an
inverse, sum over complete noncrossing pairings of its occurrences of
$M$. Remove each innermost pair by
$MDM\mapsto j n^{-1}\Tr(D)I$ and combine adjacent diagonals. An odd number
of copies of $M$ gives zero prediction; a word with no $M$ predicts itself.
For an inverse factor, apply this rule coefficientwise to
$(I-zAM+z^2j\chi_AA)^{-1}$ expanded formally at $z=0$, and then set $z=1$.
The cited lemma shows that the resulting prediction is a polynomial in $z$.
The required bounds are
\begin{equation}\label{eq:stat-word-diagnostics}
 \norm W\le C_{\rm w},\qquad
 \left(\sum_{a\ne b}|W_{ab}|^4\right)^{1/4}\le C_{\rm w},\qquad
 \norm{\diag(W-\mathsf P(W))}_2\le C_{\rm w}.
\end{equation}
The entrywise and diagonal norms are unnormalized. These bounds hold
simultaneously in all the allowed diagonal choices; $C_{\rm w}$ is
independent of $n$. We refer to them, together with $\norm M\le K$, as
the deterministic diagnostics. The construction below chooses a sufficiently
large finite $L$ and fixed coefficient bounds before $n$ tends to infinity.
Constants in the mean estimate may depend on these bounds and the fixed
stability and inverse margins.

\begin{lemma}[Arbitrary-law mean bound]\label{lem:stat-mean}
Suppose $P(x)$ is proportional to
$\exp(x^{\mathsf T}Mx/2+h^{\mathsf T}x)$ and satisfies the deterministic
diagnostics and \Cref{eq:stat-stability}.  There is a constant $C$, uniform
over these fields, matrices, and dimensions, such that every probability
$Q$ on the cube satisfies
\begin{equation}\label{eq:stat-mean}
 \norm{\E_Qx-\E_Px}\le C\bigl(1+\sqrt{D_P(Q)}\bigr).
\end{equation}
The required diagnostic length is finite and depends only on the fixed
stability margins and the bounds in the construction.
\end{lemma}

\begin{proof}
We first extend the weak residual identity to expectations under $Q$,
then use first moments to select two consecutive small residuals, and
finally compare the selected iterate with the stable root.  The
deterministic differentiation bounds are supplied by
\cite[Lemmas~5.7, 6.1, 6.3, and~6.4]{SKGap}.  The conclusion
\Cref{eq:stat-mean}, including its square-root dependence on the energy,
is the additional assertion proved in this argument.

\emph{Residual identities under the changed law.}

The recursion uses Onsager corrections, as in Bolthausen's TAP
iteration~\cite{Bolthausen2014}. Post-iteration Gaussian comparison
and local convexity are studied by Celentano~\cite{Celentano2024}; neither
result substitutes for the buffered residual estimates used here.
Define the primary fields by
\[
 \begin{gathered}
 m^0=x,\quad b_0=0,\quad h^1=h+Mx,\quad
 m^l=\tanh h^l,\\
 b_l=n^{-1}\sum_i(1-(m_i^l)^2),\qquad
 h^{l+1}=h+Mm^l-jb_lm^{l-1},
 \end{gathered}
\]
and put $R_l=m^{l-1}-m^l$.  At any fixed depth the derivative matrices of
$h^l,m^l$ have bounded operator norm and
$\norm{db_l}\le C_l n^{-1/2}$.  The small-vector constructions in \cite[Definition~6.2]{SKGap} start with
seeds of bounded Euclidean norm that may depend on $M,h$ but are
independent of the spin $x$.
Their sources are finite sums of smooth site coefficients times seeds
or earlier auxiliary fields, and are linear in those auxiliary arguments.
The coefficients and their derivatives
through every fixed order used in the construction are uniformly bounded.
Scalar parameters $\theta$ satisfy
$\sup_x(|\theta(x)|+\norm{d\theta(x)})\le C$; they may include the averages
$b_l$. A source $U$ is admissible from level $p$ when every explicit
primary field in its site coefficients has index at least $p$.
In particular, the seed $(1,\ldots,1)/\sqrt n$ permits the source
$m^l/\sqrt n$ at level $l$.
An ordinary construction starts with an empty auxiliary list.
At stage $r$, write $w^r=U$ for the new source and $y^r$ for the
output of the operation below; the auxiliary sum runs over earlier
pairs $(w^a,y^a)$, $a<r$.  The formal partial $\partial_lU_i$
differentiates the explicit primary-field argument $h_i^l$, holding
the seeds, scalar parameters (including the averages $b_l$), and
auxiliary arguments fixed.  The partial $\partial_{y^a}U_i$
differentiates the linear argument $y_i^a$ with its coefficient
held fixed.  A partial in an absent argument is zero.
The associated auxiliary operation is
\begin{equation}\label{eq:stat-auxiliary}
 U\longmapsto MU-j\sum_l\alpha_lm^{l-1}-j\sum_a\gamma_a w^a,
 \quad
 \alpha_l=n^{-1}\sum_i\partial_lU_i,\quad
 \gamma_a=n^{-1}\sum_i\partial_{y^a}U_i.
\end{equation}
These partials act on the displayed arguments of a recipe.  They are
distinct from the discrete derivative matrix $(\nabla U)_{ij}=d_jU_i$,
which records the response of the whole recipe to spin flips.
For every fixed construction the imported deterministic bounds are
\begin{equation}\label{eq:stat-small-bounds}
 \norm U+\norm{\nabla U}_{\HS}+\sum_i|d_iU_i|\le C,
 \quad
 |\alpha_l|+\norm{d\alpha_l}\le Cn^{-1/2},
 \quad |\gamma_a|+\norm{d\gamma_a}\le C.
\end{equation}
They also hold for the local implicit construction used below, on its
specified fixed flip buffers.  These statements concern pointwise
differences and have no dependence on the law under which they are tested.

Put $D=D_P(Q)$ and, for a scalar test, let
\[
 B(G)=\sup_x|G(x)|+\sup_x\norm{dG(x)}.
\]
We claim that the weak residual bound becomes
\begin{equation}\label{eq:stat-Q-residual}
 |\E_Q G R_p^{\mathsf T}U|
       \le C B(G)(1+\sqrt D).
\end{equation}
For a local construction the test is supported on its small-residual
region, and the supremum bounds are only needed on the corresponding
buffer.  In particular no conditional regularity of $Q$ is assumed.

At the first level, the only change from the equilibrium residual
argument is the mismatch between the $Q$-conditional mean and the
$P$-conditional mean.  Conditional covariance under $Q$ gives
\begin{align*}
 \E_Q\sum_i(x_i-m_i^1)GU_i
 ={}&\E_Q\sum_i(1-\bar m_i^2)d_i(GU_i)\\
 &+\E_Q G\sum_i(\bar m_i-m_i^1)U_i.
\end{align*}
This remains true when a conditional law of $Q$ is a point mass.  The
product rule $d_i(GU_i)=Gd_iU_i+(d_iG)(U_i-2x_id_iU_i)$ and
\Cref{eq:stat-small-bounds} bound the first term by $CB(G)$.
\Cref{eq:stat-conditional-means} and Cauchy--Schwarz bound the second by
$2C\norm G_\infty\sqrt D$.  This proves the base case.

For higher levels we reduce to lower residual indices by an exact
cancellation.  For a target admissible from level $p+1$, set
\[
 a_p=j(b_p-b_{p-1}),\qquad
 U'_i=\Phi(h_i^p;h_i^{p+1},a_p)U_i,
\]
where
\begin{equation}\label{eq:stat-Phi}
 \Phi(z;r,a)=\int_0^1 e^{a(1-u^2)/2}
       \frac{\cosh(u(z-r))}{\cosh z\cosh r}\,\dd u.
\end{equation}
It satisfies
$[(z-r-a\tanh z)-a\partial_z]\Phi=\tanh z-\tanh r$.
Adjoin the auxiliary field $y'$ of $U'$ using
\Cref{eq:stat-auxiliary}.  With $\alpha_l,\gamma_a$ belonging to $U'$,
the recurrence gives
\begin{align}
 R_{p+1}^{\mathsf T}U
 ={}&R_p^{\mathsf T}y'+a_pR_p^{\mathsf T}U'
   -jb_{p-1}R_{p-1}^{\mathsf T}U'
   +j\sum_a\gamma_a R_p^{\mathsf T}w^a\notag\\
 &+j\sum_{l\ge p}\alpha_lR_p^{\mathsf T}m^{l-1}
   -a_p n\alpha_p.\label{eq:stat-residual-induction}
\end{align}
The $R_{p-1}$ term is absent at $p=1$.  The $l=p$ term cancels the last
term because
\[
 R_p^{\mathsf T}m^{p-1}
       =n(b_p-b_{p-1})-R_p^{\mathsf T}m^p.
\]
After this cancellation, the remaining $l=p$ term is
$-j(\sqrt n\alpha_p)R_p^{\mathsf T}(m^p/\sqrt n)$.
Its target $m^p/\sqrt n$ is admissible from level $p$ and has bounded
norm. For $l>p$, use the target $m^{l-1}/\sqrt n$ in the same way.
Thus every term has a lower residual index, a permitted target, and one
of the bounded multipliers $a_p,b_{p-1},\gamma_a,\sqrt n\alpha_l$.
Their difference bounds in \Cref{eq:stat-small-bounds} and the product
rule give $B(\theta G)\le C B(G)$ for each such multiplier.
Strong induction over finite admissible constructions therefore proves
\Cref{eq:stat-Q-residual}.  Each step adds finitely many operations and
decreases the residual index, so a finite word length suffices.  For a
supported local test, a difference can be nonzero only on a one-flip
neighborhood of its support; the imported buffers include every such
neighborhood used in this finite induction.

\emph{Selecting two consecutive small residuals.}
The weak estimate now has the desired $1+\sqrt D$ cost.  We preserve that
cost in the selection step by using first moments rather than second
moments.  Set $S_l=R_l^{\mathsf T}m^l/\sqrt n$ and
$\psi(s)=s/\sqrt{1+s^2}$.  The primary derivative bounds give
$\sup\norm{dS_l}\le C_l$, while $|\psi|,|\psi'|\le1$ and
$s\psi(s)\ge|s|-1$.  Testing \Cref{eq:stat-Q-residual} with
$U=m^l/\sqrt n$ and $G=\psi(S_l)$ yields
\begin{equation}\label{eq:stat-S-first-moment}
 \E_Q|S_l|\le C_l(1+\sqrt D).
\end{equation}
The exact telescoping identity is
\[
 \norm{R_l}^2/n=b_l-b_{l-1}-2S_l/\sqrt n.
\]
For a fixed small $\rho>0$, choose a fixed depth $d_0$ with
$\lfloor d_0/2\rfloor\rho^2/4>2$, and then a fixed $\varepsilon>0$
with $2d_0\varepsilon<1$.  If every $|S_l|\le\varepsilon\sqrt n$,
the sum of the squared residuals divided by $n$ is less than two.
Some consecutive pair consequently has both norms at most
$\rho\sqrt n/2$.  Let $C_k$ be smooth cutoffs equal to one on this latter
condition and supported where both norms are at most $\rho\sqrt n$.
The primary derivative bounds give $\norm{dC_k}\le C_\rho/\sqrt n$.
Combine the cutoffs in order so that earlier cutoffs reduce the weight
available to later ones, and retain the unassigned deficit.  Put
\[
 \widehat C_k=C_k\prod_{l=2}^{k-1}(1-C_l),\qquad
 \Delta=1-\sum_{k=2}^{d_0}\widehat C_k.
\]
Markov's inequality and \Cref{eq:stat-S-first-moment} give
\begin{equation}\label{eq:stat-selection}
 Q(\Delta>0)\le C_\rho(1+\sqrt D)/\sqrt n.
\end{equation}

\emph{Comparing the selected iterate with the stable root.}
Let $r$ be the unique zero of \Cref{eq:stat-field-map}, whose existence
and inverse bound follow from \cite[Lemma~5.7]{SKGap}, and put
$t_* =\tanh r$, $q_*=n^{-1}\norm{t_*}^2$.  On the buffered region of
$C_k$, fix the primary fields and the root, and choose one of the two
sources $w_0$ specified below. The implicit construction of
\cite[Lemma~6.3]{SKGap} supplies a vector $w$, an auxiliary vector
$y^{\rm imp}$, and a scalar $c_1$ solving
\[
 \begin{gathered}
 a=j(1-b_{k-1}-q_*),\quad A_i=\Phi(h_i^k;r_i,a),\quad
 A=\diag(A_i),\quad \chi=n^{-1}\sum_iA_i,\\
 w=w_0+A(y^{\rm imp}-jc_1t_*),\qquad
 y^{\rm imp}=Mw-j\chi w-jc_1m^{k-1},\qquad
 c_1=n^{-1}\sum_i\partial_k w_i.
 \end{gathered}
\]
Here $w_0$ is either $Ae$ for a fixed $\norm e\le1$, or
$n^{-1/2}\partial_r\Phi(h_i^k;r_i,a)$ coordinatewise. In
$\partial_k w_i$, the arguments $y_i^{\rm imp},c_1,a$ are held fixed;
in $\partial_r\Phi$, $a$ is held fixed. With this convention the displayed
equations form a linear system for $(w,y^{\rm imp},c_1)$, whose unique
solution is supplied by the cited lemma. That lemma also gives
\[
 \norm w+\norm{y^{\rm imp}}+\sqrt n|c_1|\le C,\qquad
 I+j\chi A-A^{1/2}MA^{1/2}\succeq c_0I
\]
after decreasing the common positive inverse margin if necessary.
The distinction between its bounds matters for selecting parameters:
derivative constants may depend on the fixed depth, while size and
inverse constants do not.  We will first control the susceptibility
average with the second seed, then test the mean in direction $e$ with
the first seed.

Define
\[
 T_k(w)=\{h^k-h-(M-jb_{k-1}I)m^k\}^{\mathsf T}w
                 -j(b_k-b_{k-1})nc_1.
\]
The algebra in \cite[Lemma~6.6]{SKGap} expresses this as
\[
 R_k^{\mathsf T}y^{\rm imp}
 +j(\chi-b_{k-1})R_k^{\mathsf T}w
 -jb_{k-1}R_{k-1}^{\mathsf T}w-jc_1R_k^{\mathsf T}m^k.
\]
Thus the already proved local residual estimate gives
\begin{equation}\label{eq:stat-Q-terminal}
 |\E_Q G T_k(w)|\le CB(G)(1+\sqrt D).
\end{equation}
This is the required arbitrary-law version of the terminal residual
identity; the last term uses the multiplier $\sqrt n c_1$.

To extract a mean comparison from this terminal residual, write
$p=m^k-t_*$ and $L_* =\sqrt n\{\chi-(1-q_*)\}$.  Subtract the
root equation and substitute the displayed equation for $Mw$ to obtain
\begin{align}
 T_k(w)={}&\sum_i[(h_i^k-r_i-am_i^k)-a\partial_k]w_{0,i}\notag\\
 &+j(1-q_*-\chi)p^{\mathsf T}w-jc_1R_k^{\mathsf T}p.
 \label{eq:stat-implicit-comparison}
\end{align}
The inverse bound and the two small residuals imply
$\norm p\le C\rho\sqrt n$.  Consequently the middle term has absolute
value at most $C_0\rho|L_*|$, with $C_0$ independent of the selection
depth.  Also $\sup\norm{dL_*}\le C_k$.
Because $C_0$ is independent of depth, we can choose $\rho$ with
$C_0\rho<1/2$ first, and only then choose $d_0$ and $\varepsilon$ as in
the selection step above.

For $w_{0,i}=n^{-1/2}\partial_r\Phi(h_i^k;r_i,a)$ the first sum in
\Cref{eq:stat-implicit-comparison} is $L_*$: differentiate the identity
following \Cref{eq:stat-Phi} in $r$, holding $a$ fixed.
Test with $G=C_k\psi(L_*)$.  Its $B$-norm is bounded.  The last term is
controlled by \Cref{eq:stat-Q-residual} with target $p/\sqrt n$ and
multiplier $\sqrt n c_1$.  Absorbing the middle term gives
\begin{equation}\label{eq:stat-L-first-moment}
 (1-C_0\rho)\E_Q[C_k|L_*|]\le1+C_k'(1+\sqrt D).
\end{equation}
For $w_0=Ae$, the first sum is $p^{\mathsf T}e$.  Test this time with
$\widehat C_k$ alone.  Since $\widehat C_k\le C_k$,
\Cref{eq:stat-Q-terminal,eq:stat-L-first-moment} and the same residual
bound imply
\[
 |\E_Q\widehat C_k(m^k-t_*)^{\mathsf T}e|
       \le C_k'(1+\sqrt D).
\]
It remains to pass from the selected iterate back to the original spin.
Sum \Cref{eq:stat-Q-residual} with constant target $e$ to handle
$x-m^k=\sum_{l\le k}R_l$, and then sum over $k$.  The omitted part has
absolute value at most
$2\sqrt n\,Q(\Delta>0)\le C(1+\sqrt D)$ by
\Cref{eq:stat-selection}.  Taking the supremum over unit $e$ proves
$\norm{\E_Qx-t_*}\le C(1+\sqrt D)$.
Apply this also with $Q=P$, when $D=0$, and use the triangle inequality.
\end{proof}

\subsection{Two observation stages and simultaneous posterior gaps}

We next construct the observation along which the mean bound can be
used.  White observation first creates a large field; a second observation
then removes the interaction while preserving quantitative stability.
Fix $k>K$, where $\norm J\le K$ on the matrix event under consideration,
and put $A_*=J+kI\succ0$.  Sample $X$ from $\mu$.  With independent
standard Brownian motions $B,W$, observe first
\[
 Y_s=sX+B_s\quad(0\le s\le b)
\]
and then, conditionally on the first observations,
\[
 Z_r=rA_*X+A_*^{1/2}W_r\quad(0\le r\le1).
\]
The posterior at the two stages has, respectively, interaction and field
\begin{equation}\label{eq:stat-posterior-parameters}
 (J,Y_s),\qquad ((1-r)J,Y_b+Z_r).
\end{equation}
Indeed the likelihood of the second observation subtracts
$r x^{\mathsf T}A_*x/2$ from the quadratic Hamiltonian; its diagonal part
is constant on the cube.  The posterior at time $b+1$ is therefore the
product measure with coordinate logits $h=Y_b+Z_1$.

Before estimating the posterior gaps, we record the two observation
identities used in their proof. The relative entropy of the observation
law is a drift energy, and the expected posterior Dirichlet energy does
not increase. In the experiment above the instantaneous precision is
$I$ on the first stage and $A_*$ on the second. More generally, let it be
$A_t\succeq0$,
with bounded operator norm, and let $P_t,Q_t$ be its posteriors for inputs
$P,Q$.  Write $\mathbb P,\mathbb Q$ for the observation laws.
The entropy-as-drift-energy identity is classical; see
Lehec~\cite[Theorem~2]{Lehec2013}, which credits F\"ollmer. Here the
innovation process is Brownian and its changed drift is the difference
of posterior means. For a bounded stopping time $\sigma$, the density
martingale and It\^o's formula give
\begin{equation}\label{eq:stat-observation-entropy}
 H(\mathbb Q|_{\mathcal F_\sigma}\mid
       \mathbb P|_{\mathcal F_\sigma})
 =\frac12\E_{\mathbb Q}\int_0^\sigma
  \norm{A_t^{1/2}(\E_{Q_t}x-\E_{P_t}x)}^2\,\dd t.
\end{equation}
For example represent the signal as
$\dd Y_t=A_tX\,\dd t+A_t^{1/2}\dd B_t$; its innovation has the same
covariance, and the density score is the difference of posterior means.
Finite dimension and bounded horizon justify optional stopping and the
change of drift.  The same formula follows for singular $A_t$ by
restricting to its range.  At deterministic times, and also at bounded stopping times,
\begin{equation}\label{eq:stat-energy-monotone}
 \E_{\mathbb Q}D_{P_t}(Q_t)\le D_P(Q).
\end{equation}
To see this, set $f=\sqrt{Q/P}$.  The observation density is $P_tf^2$,
so the left side is $\E_{\mathbb P}\mathcal E_{P_t}(f)$.
Its $i$th summand is
$\E\Var(f(X)\mid X_{-i},\mathcal F_t)$, bounded by
$\E\Var(f(X)\mid X_{-i})$ by conditional variance decomposition. The same
argument applies with $\mathcal F_t$ replaced by the stopped observation
$\sigma$-field.
These arguments are the finite-state identities of
\cite[Section~2]{SKGap}, with a bounded precision inserted.

\begin{proposition}[Observation paths and cavities]\label{prop:stat-posterior-gap}
There are fixed $b,k,C,c,\gamma>0$ and events of disorder probability
tending to one on which the following assertions hold, for the full model
and for every model obtained by deleting one site.  Along the observation
experiment just defined, all posterior heat-bath gaps are at least
$\gamma$, except on an ordinary observation event of probability at most
$Ce^{-cn}$.  On the same good paths all posteriors satisfy
\Cref{lem:stat-mean} with the same constant.  The exceptional event and
a stopping time before loss of these properties can be chosen independently
of any test function.  The constants are common to all cavities.
\end{proposition}

\begin{proof}
We verify second-stage stability and its gap first.  We then run the
observation argument conditionally from each first-stage posterior to
obtain its gap, and finally place all single-site cavities on one common
disorder event.

\emph{Stability through the two stages.}
For the first stage, \cite[Proposition~5.1 and Corollary~5.2]{SKGap} give
\Cref{eq:stat-stability}, with positive slack and a test-independent stopping
rule, throughout the first stage up to any prescribed fixed $b$, outside
an event of ordinary observation probability $e^{-cn}$.  Intersect its
matrix event with the finite word events required in
\Cref{lem:stat-mean} and in the directional estimate
\cite[Proposition~7.2]{SKGap}.

On the second stage write $u=1-r$ and use $M=uJ$, $j=u^2\beta^2$.
The word bounds hold uniformly in $u\in[0,1]$.  In fact their inverse
factor is exactly
\[
 (I-uAJ+u^2\beta^2\chi_AA)^{-1}
   =(I-(uA)J+\beta^2\chi_{uA}(uA))^{-1}.
\]
The other word factors only acquire powers of $u$.  The fixed upper bound
on $A$ can thus be the same as for $J$.  All derivative and implicit
construction estimates used above are uniform for these parameters:
the recurrences have bounded coefficients, $j$ remains bounded away from
one, and the inverse margins are fixed.  At $u=0$ the measure is product;
\Cref{eq:stat-conditional-means} gives its mean bound directly.

To obtain a margin throughout the second stage, use the large field
created in the first one.  Gaussian
maximal bounds on the two fixed Brownian intervals imply, outside
$Ce^{-cn}$, that
\begin{equation}\label{eq:stat-large-field}
 \sup_{0\le r\le1}\norm{Y_b+Z_r-bX}\le C\sqrt{bn}
\end{equation}
for $b\ge1$, with a constant depending on $K,k$.  One elementary proof of
the bound for $B$ uses a $1/2$-net of the unit sphere and the scalar
Brownian maximal inequality; the covariance norm of $A_*^{1/2}W$ is at
most $K+k$.  This also gives the asserted exponential probability.
If $y$ has residual at most $\rho_0\sqrt n$ in
\Cref{eq:stat-field-map}, then
$\norm{y-h}\le(\rho_0+\beta^2+K)\sqrt n$.
Thus, by \Cref{eq:stat-large-field}, at most $C n/b$ coordinates have
$|y_i|<b/2$.  Outside these coordinates $V(y)_{ii}\le\sech^2(b/2)$.
For a candidate $A$ in \Cref{eq:stat-stability}, an additional set of size
at most $\epsilon_A^2n/a^2$ has $|A_{ii}-V(y)_{ii}|>a$.

The small principal submatrix bounds of \cite[Lemma~8.2]{SKGap} apply also
to $J$, because $\norm{J-(\tanh J_{ij})}\to0$ on that event.  If
$S$ is the preceding exceptional set and $A_{ii}\le2a$ off $S$, they give
\begin{align*}
 \norm{A^{1/2}uJA^{1/2}}
 &\le a_+\norm{J_{SS}}+2K\sqrt{2a_+a}+2Ka,\\
 \frac{2u^2\beta^2}{n}\norm{A^{1/2}\tanh y}^2
 &\le2\beta^2\{a_+|S|/n+2a\}.
\end{align*}
Choose the principal-submatrix tolerance and $a$ first so that the sum of
these bounds is less than $1/2$; choose an allowed exceptional fraction
next, then $b$ large and $\epsilon_A$ small.  The nonnegative term
$jb(y)A$ in \Cref{eq:stat-field-hessian} can be kept or dropped.
This proves \Cref{eq:stat-stability} with $\kappa=1/2$ on the second
stage, after reducing the common margins if necessary.  These choices
are compatible with the first-stage theorem: choose this fixed $b$ first
and subsequently intersect and reduce its supplied margins.

\emph{The second-stage gap.}
The same large-field event also gives a gap at least $1/2$ throughout
this stage.  The curvature bound must hold uniformly in the spin
configuration.  For that purpose combine \Cref{eq:stat-large-field} with
$\norm{uJx}\le K\sqrt n$.  Only a small fixed fraction of the fields
$h_i+(uJx)_i$ can be smaller than $b/4$ in absolute value.
Let $v_i=\sech^2(h_i+(uJx)_i)$ and $a_i=v_id_if$.
The signed curvature inequality of \cite[Lemma~8.1]{SKGap} is
\[
 \E_P(Hf)^2\ge(1-Cr_4)\mathcal E_P(f)
       -\E_P\{a^{\mathsf T}Ua+C|a|^{\mathsf T}Q_0|a|\},
\]
where $H=-L_P$, $U_{ij}=\tanh(uJ_{ij})$,
$(Q_0)_{ij}=U_{ij}^2$, and $r_4=\max_i\sum_jU_{ij}^4=o(1)$.
On the weak-field set use the small principal-submatrix norms; on its
complement use $v_i\le\sech^2(b/4)$.  If
$q_f(x)=\sum_i v_i(d_if)^2$, each quadratic cost is bounded by
$[\varepsilon+K'(2\sqrt\eta+\eta)]q_f(x)$ with
$\eta=\sech^2(b/4)$ and a fixed $K'$.  Choosing the same tolerances
sufficiently small proves
$\E_P(Hf)^2\ge\mathcal E_P(f)/2$.  The finite-dimensional spectral
theorem gives the gap $1/2$.  The matrix bounds are uniform in $u$:
$\tanh(uJ)=uJ+o_{\op}(1)$, while the nonnegative matrix
$(\tanh(uJ_{ij})^2)$ is dominated entrywise by $(J_{ij}^2)$.

\emph{The first-stage posterior gaps.}
The remaining task is to obtain a gap at every first-stage posterior.
We do this by conditioning the observation experiment from that posterior
onward.  The zero-field gap theorem does not by itself apply to these
random fields.  Let $E$ be the event of a first-stage stability
failure or failure of the terminal gap at time $b$.  The preceding bounds
and \cite[Proposition~8.3]{SKGap} allow a fixed $a_0>0$ such that
$\P(E\mid J)\le e^{-2a_0 n}$.  For the observation filtration set
$q_s=\P(E\mid\mathcal F_s,J)$.  The nonnegative martingale maximal
inequality gives
\begin{equation}\label{eq:stat-future-bad}
 \P\{\sup_{s\le b}q_s>e^{-a_0n}\mid J\}\le e^{-a_0n}.
\end{equation}
Outside this event and $E$, fix a time $s<b$.  Conditional on
$\mathcal F_s$, the remaining experiment is an ordinary white
observation starting from its posterior $P_s$, with horizon $b-s\le b$.
Its future bad probability is at most $e^{-a_0n}$.

For every $f$ with $P_sf=0$ and $P_sf^2=1$, set
$d=\mathcal E_{P_s}(f)$, and stop the remaining experiment at its next
stability failure or at $b$, denoting this time by $\sigma$.  Write
$N_t=P_tf^2$ and $V_t=P_tf^2-(P_tf)^2$.
Conditional variance decomposition gives
$\E[\mathcal E_{P_t}(f)\mid\mathcal F_s]\le d$.
At every retained stable posterior, the directional estimate of
\cite[Proposition~7.2]{SKGap} states, in squared form,
\[
 \norm{\Cov_{P_t}(f,x)}^2
 \le C_{\rm d}\{\Var_{P_t}(f)+\mathcal E_{P_t}(f)\}.
\]
It uses the same stable-field implication and a sufficiently long finite
set of word diagnostics. This homogeneous estimate, combined with the
filtering variance identity, implies
\begin{equation}\label{eq:stat-conditional-variance}
 \E[V_\sigma\mid\mathcal F_s]
 \ge v_s-(1-v_s)d,
 \qquad v_s=e^{-C_{\rm d}(b-s)}\ge v_0=e^{-C_{\rm d}b}>0.
\end{equation}
Indeed the stopped expected variance has derivative
$-\E[\1_{t<\sigma}\norm{\Cov_{P_t}(f,x)}^2\mid\mathcal F_s]$,
which is at least $-C_{\rm d}(\E V_{t\wedge\sigma}+d)$.

The variance lower bound must be compared with the terminal energy.
To control the variance lost on bad paths, change the input law to
$f^2P_s$.  The observation density is $N_t$
and its posterior is $Q_t=f^2P_t/N_t$.  Before stopping,
\Cref{lem:stat-mean} and conditional variance decomposition give
\[
 \E^{f}[\1_{t<\sigma}\norm{\E_{Q_t}x-\E_{P_t}x}^2]\le C(1+d),
 \qquad
 \E^{f}D_{P_t}(Q_t)
 =\E\mathcal E_{P_t}(|f|)\le d.
\]
By \Cref{eq:stat-observation-entropy}, the stopped observation entropy is
at most $Cb(1+d)$.  Relative entropy between event probabilities gives
\begin{equation}\label{eq:stat-changed-bad}
 \P^{f}(\mathrm{future\ bad}\mid\mathcal F_s)
 \le\frac{Cb(1+d)+\log2}{a_0n}.
\end{equation}
Here the bad event is measurable at $\sigma$: it is the stopping event
before $b$, or the terminal gap failure at $b$.  Its ordinary conditional
probability is exponentially small by \Cref{eq:stat-future-bad}.
The changed probability in \Cref{eq:stat-changed-bad} is derived anew for
each $f$; it is not asserted to be exponentially small.

On good terminal paths $V_b\le2\mathcal E_{P_b}(f)$; on bad paths
$V_\sigma\le N_\sigma$.  The stopped density identity and
\Cref{eq:stat-conditional-variance,eq:stat-changed-bad} consequently give,
for $d\le1$,
\[
 v_0-(1-v_0)d\le2d+C/n.
\]
For all sufficiently large $n$ this implies $d\ge v_0/6$; for $d>1$
the same lower bound is automatic.  All constants are uniform in $s$ and
$f$.  A short remaining horizon improves $v_s$ and decreases the entropy
integrals, so the estimate stays uniform as $s\uparrow b$.  At $s=b$ use
the terminal gap directly.  Alternatively one can first argue at rational
times and then use continuity of the finite-state gap in the field.

\emph{One event for every cavity.}
We now transfer the matrix and observation estimates to all single-site
cavities at once.  This requires a common disorder event, rather than a
union bound over unspecified $o(1)$ disorder probabilities.  Write
$J^{(i)}$ for a principal minor.  Word diagnostics for this matrix are restrictions of those for
$J$: insert the complementary coordinate projection among diagonal factors,
and set the missing entry of an inverse diagonal to zero.  All normalized
traces can be retained with denominator $n$; equivalently the $(n-1)$-site
variance parameter is $j_i=\beta^2(n-1)/n$.  The identity
$j_i/(n-1)=\beta^2/n$ makes these the same contractions.
Inserting these projections may increase the required word length by a
fixed factor; choose that larger finite length in the full event. The
event is already simultaneous in its diagonal parameters.  The small-submatrix bounds pass directly to
every minor.

For first-stage stability, take a candidate cavity vector $y_{-i}$ and
extend it by choosing $y_i$ to solve the missing scalar equation
$F_{J,h}(y)_i=0$.  This equation is continuous with a bounded nonlinear
part and hence has a solution.  On the retained coordinates the residual
changes by a vector of norm at most
$\norm{J_{-i,i}}+O(n^{-1/2})=O(1)$.
Extend the candidate diagonal with $A_{ii}=\sech^2 y_i$.
Its Frobenius discrepancy is unchanged.  The principal block of the full
Hessian differs from the cavity Hessian by $O(1/n)$ in norm: the rank
coefficient is exactly $2\beta^2/n$ in either dimension, and only the
missing contribution to $jb(y)$ changes.  Start with smaller cavity
residual and Frobenius thresholds and use the slack in the full stability
event.  It follows, uniformly in $i$, that full stability implies cavity
stability with common positive margins for large $n$.

The comparison of sampling laws costs only $\exp(C\sqrt n)$.  Indeed the
full marginal on $x_{-i}$ is proportional to
$\mu^{(i)}(x_{-i})\cosh(J_i x_{-i})$, and
$|J_i x_{-i}|\le K\sqrt n$.  More explicitly, the law obtained by sampling
the cavity spin configuration and an independent uniform missing spin has
density at most $\exp(2K\sqrt n)$ relative to $\mu$.  Extend its Brownian
observation by an independent coordinate.  Thus a full exponential
observation failure remains exponential for the cavity experiment.
The second-stage large-field and curvature arguments apply directly to
the minors with the same constants.  A union bound over the $n$ cavities
now loses only a factor $n$ from exponential observation probabilities.
This proves the claimed simultaneous statement on one common disorder
event.  It also proves a common equilibrium gap by taking $s=0$.

\emph{A test-independent stopping rule.}
To finish, all properties tested above depend only on the observed field
and finite matrices.  Stable-field failure is detected by the Lipschitz
residual-distance stopping construction in
\cite[Corollary~5.2]{SKGap}; the gap is a continuous finite-matrix
eigenvalue.  Stopping at smaller stability margins or at a gap below half
the established bound gives an observation stopping time independent of
the test function, with the same exponential exceptional bound.
\end{proof}

\subsection{Entropy through the product endpoint}

We now combine the observation identities and posterior mean bound with
entropy tensorization at the product endpoint.

At the product endpoint, the one-site entropy constants grow with the
coordinate logits.  We therefore need a bound on sums of their squares
over every subset of sites.  Its good-disorder event may depend on a fixed
accuracy $\delta$, unlike the posterior-gap event.

\begin{lemma}[Subset tails for the terminal field]\label{lem:stat-field-tails}
There is $C<\infty$, independent of $\delta>0$, such that, on events of
disorder probability tending to one for each fixed $\delta$, the ordinary
two-stage experiment satisfies
\begin{equation}\label{eq:stat-field-tails}
 \mathbb P\left\{\exists I\subset[n]:
  \sum_{i\in I}h_i^2>C\left[|I|\log\frac{en}{|I|}+s\right]\right\}
       \le e^{-s}\qquad(s\ge\delta n).
\end{equation}
The empty subset contributes zero.  The same event and constants can be
taken for all single-site cavities, with their own dimension in the formula.
\end{lemma}

\begin{proof}
We prove a Gaussian subset bound under the planted law, transfer it to
stationary sampling using the partition function, and then convert the
averaged bound to one at fixed disorder.  The exponential slack will also
make the cavity statement simultaneous.  Let $Z(J)$
be the partition function with counting measure on the cube.  The annealed
planted law has density
$e^{x^{\mathsf T}Jx/2}/\E_J Z(J)$ relative to Gaussian disorder times
counting measure on $x$.  Under it $X$ is uniform, and, conditionally on
$X$, the off-diagonal entries are independent Gaussian variables with mean
$\beta^2X_iX_j/n$ and variance $\beta^2/n$.  Moreover
\begin{equation}\label{eq:stat-Z-moment}
 \frac{\E_J Z(J)^2}{(\E_J Z(J))^2}
 =e^{-\beta^2/2}\E\exp\left(\frac{\beta^2 R_n^2}{2n}\right)
 \le(1-\beta^2)^{-1/2},
\end{equation}
where $R_n$ is the sum of $n$ independent uniform signs.  The last bound
follows by introducing a standard Gaussian $G$ and using
$\cosh t\le e^{t^2/2}$:
$\E e^{\beta^2R_n^2/(2n)}=\E_G\cosh(\beta G/\sqrt n)^n$.
The gradient of $\log Z$ in the independent standardized Gaussian
coordinates has norm at most $C\sqrt n$.  Gaussian concentration and
Paley--Zygmund applied to \Cref{eq:stat-Z-moment} imply
\begin{equation}\label{eq:stat-Z-comparison}
 \log Z(J)-\log\E_JZ(J)=o(n)
       \quad\hbox{in disorder probability}.
\end{equation}
For clarity, Paley--Zygmund gives a fixed positive probability of
$\log Z\ge\log\E Z-\log2$; concentration places its mean within
$O(\sqrt n)$ of this level.  Concentration once more proves
\Cref{eq:stat-Z-comparison}.

We next estimate the fields under planting.  At the observation endpoint
\[
 h=(b+k)X+JX+B_b+A_*^{1/2}W_1.
\]
Under planting, conditionally on $X$, $JX$ is a Gaussian vector with
bounded covariance norm and a mean of bounded coordinate magnitude.
On $\norm J\le K$, the last observation noise has conditional covariance
norm at most $K+k$.  Thus, for a fixed subset of size $m$, Gaussian norm
tails, followed by a union bound over at most $(en/m)^m$ subsets, give
an annealed bound $C_1e^{-c_1t}$ for violation of
\[
 \sum_{i\in I}h_i^2\le C_2[m\log(en/m)+t]
\]
simultaneously in $I$.  This uses the deterministic inequality
$\norm{z_1+\cdots+z_4}^2\le4\sum\norm{z_a}^2$, so independence of
$JX$ and $A_*^{1/2}W_1$ is unnecessary.  Increasing $C_2$ gives any
fixed desired exponential rate in $t$. Fix this increase once, before the
quenched comparison below; it does not depend on $\delta$.

To return to the quenched experiment, note that the density of the
planted joint spin-disorder law relative to stationary sampling with the
original disorder is $Z(J)/\E_JZ(J)$.  Restrict to
\Cref{eq:stat-Z-comparison} with, for example,
$\log Z\ge\log\E Z-\delta n/10$.  The preceding annealed bound, with
ample exponential slack, therefore bounds the disorder average of the
quenched failure probability at level $s$ by $e^{-3s}$ for
$s\ge\delta n$, once $n$ is large.  Markov's inequality and a union bound
over the integer levels $s\ge\delta n$ give a disorder exceptional
probability $O(e^{-c\delta n})$ outside which the required conditional
probability is at most $e^{-s}$.  For a real $s\ge\delta n\ge1$, apply
the integer estimate at $\lceil s\rceil$ and double the threshold
constant; then $\lceil s\rceil\le2s$ and
$e^{-\lceil s\rceil}\le e^{-s}$.
Only the good disorder event depends on $\delta$.

For a cavity the same planted computation has parameter
$\beta\sqrt{(n-1)/n}<\beta$ and common constants.  Alternatively extend
its spin and observation experiment as in
\Cref{prop:stat-posterior-gap}; the additional
$\exp(C\sqrt n)$ factor is absorbed by the exponential slack for fixed
$\delta$.  Taking the union over cavities still gives probability tending
to one.  This argument avoids a union bound over unquantified convergence
in probability in \Cref{eq:stat-Z-comparison}: the full partition lower
bound transfers to the minors with an $O(\sqrt n)$ error.  More explicitly,
\[
 2Z(J^{(i)})\le Z(J)\le2e^{K\sqrt n}Z(J^{(i)}),\qquad
 \log\frac{\E Z(J)}{2\E Z(J^{(i)})}
       =\frac{\beta^2(n-1)}{2n}.
\]
These inequalities transfer the same full lower bound to every minor at
once; the remaining exponential Markov bounds tolerate the union over $i$.
\end{proof}

We can now combine the drift-energy identity with tensorization at the
product endpoint.  The subset estimate controls the field weights without
paying for the largest coordinate logit.

\begin{theorem}[Entropy-energy estimate]\label{thm:stat-entropy}
There is a constant $C$ such that for each fixed $\delta>0$, on events of
disorder probability tending to one, every probability $\nu$ satisfies
\begin{equation}\label{eq:stat-entropy}
 H_\mu(\nu)\le C\left\{\delta n+1+
 D_\mu(\nu)\left(1+\sqrt{\log\frac{en}{D_\mu(\nu)}}\right)\right\}.
\end{equation}
The product containing $D_\mu(\nu)$ is zero when this energy is zero.
The constant $C$ is independent of $\delta$ and $\nu$.  The assertion
holds on one common event for all single-site cavities as well.
\end{theorem}

\begin{proof}
We first bound the entropy of the observation law, then the conditional
entropy left at its product endpoint.  The latter bound contains a term
involving the original entropy, which will be absorbed only at the end.

\emph{Entropy of the observations.}
Put $D=D_\mu(\nu)$ and $H=H_\mu(\nu)$, and change the input of the
two-stage experiment from $\mu$ to $\nu$.  Let $P_t,Q_t$ denote the
ordinary and changed posteriors, and let $\sigma$ be the first loss of
the retained stability conditions, capped at $b+1$.
By \Cref{lem:stat-mean,eq:stat-observation-entropy,eq:stat-energy-monotone},
\begin{equation}\label{eq:stat-stopped-H}
 H(\mathbb Q|_{\mathcal F_\sigma}\mid
       \mathbb P|_{\mathcal F_\sigma})\le C(1+D).
\end{equation}
The ordinary probability of $\sigma<b+1$ is at most $Ce^{-cn}$.
Applying the binary entropy inequality to this stopped event gives
\[
 \mathbb Q\{\sigma<b+1\}\le C(1+D)/n.
\]
Continue the observations even on this event.  The squared difference of
posterior means is always at most $4n$, and the precision and horizon are
bounded.  Hence the contribution after stopping in
\Cref{eq:stat-observation-entropy} is at most another $C(1+D)$.
The relative entropy of the complete observation laws is therefore at most
$C(1+D)$.

\emph{Entropy remaining at the product endpoint.}
The observation channel is the same for the two input laws.  Thus the
entropy chain rule on the joint experiment reads
\begin{equation}\label{eq:stat-entropy-chain}
 H=H(\mathbb Q\mid\mathbb P)
       +\E_{\mathbb Q}H_{P_{b+1}}(Q_{b+1}).
\end{equation}
At the product endpoint, entropy tensorization gives
\[
 H_{P_{b+1}}(Q_{b+1})
 \le\sum_i\E_{(Q_{b+1})_{-i}}
          H_{(P_{b+1})_i}((Q_{b+1})_i).
\]
This version with all other coordinates conditioned follows, for example,
by the entropy chain rule and convexity of relative entropy: conditioning
on more coordinates can only increase the corresponding average conditional
entropy relative to a fixed coordinate law.

For a two-point law $p$ of least atom $p_{\min}$ and any $q$,
\begin{equation}\label{eq:stat-two-point}
 H_p(q)\le C(1+\log(1/p_{\min}))
       \left[1-\left(\sum\sqrt{pq}\right)^2\right].
\end{equation}
Indeed for $0\le z\le1/p_{\min}$,
$z\log z-z+1\le C(1+\log(1/p_{\min}))(\sqrt z-1)^2$:
the ratio has a finite limit at $z=1$, is bounded on $[0,4]$, and for
$z\ge4$ is at most $C\log z$.  Sum this inequality with weights $p$ and
use $\sum p(\sqrt{q/p}-1)^2=2(1-\sum\sqrt{pq})
\le2[1-(\sum\sqrt{pq})^2]$.
For the product law of logit $h_i$,
$\log(1/p_{\min})\le\log2+2|h_i|$.
Define the conditional energy contributions
\[
 z_i=\E_{(Q_{b+1})_{-i}}
  \left[1-\left(\sum_{s=\pm1}
   \sqrt{(P_{b+1})_i(s)(Q_{b+1})_i(s)}\right)^2\right].
\]
They satisfy $0\le z_i\le1$ and
$\E_{\mathbb Q}\sum_i z_i\le D$ by
\Cref{eq:stat-energy-monotone}.  Consequently
\begin{equation}\label{eq:stat-H-z}
 H\le C(1+D)+C\E_{\mathbb Q}\sum_i z_i|h_i|.
\end{equation}

\emph{The weighted field sum and absorption.}
The variables $z_i$ record where the energy is spent.  We use subset
tails to bound their field-weighted sum without replacing the field by
its maximum over sites.  By \Cref{lem:stat-field-tails} there is a
nonnegative excess $R(h)$,
at least $\delta n$, such that simultaneously in all subsets
\[
 \sum_{i\in I}h_i^2\le C\left[|I|\log\frac{en}{|I|}+R(h)\right],
 \qquad \mathbb P\{R(h)>\delta n+t\}\le e^{-t}\quad(t\ge0),
\]
after a fixed increase of $C$.  In particular
$\E_{\mathbb P}e^{R/2}\le C e^{\delta n/2}$.
The entropy variational inequality on the joint spin-observation laws,
whose relative entropy is exactly $H$, gives
\begin{equation}\label{eq:stat-excess}
 \E_{\mathbb Q}R\le C(\delta n+1+H).
\end{equation}
The same conclusion follows from the observation marginal by data
processing.  At this stage $H$ is still unknown.  Keeping it on the right
side is legitimate; the square-root estimate below will allow us to absorb the
resulting contribution.

Set $z=\sum_i z_i$.  Layer cake for $I_t=\{i:z_i>t\}$ and concavity of
$u\mapsto u\log(en/u)$ on $[0,n]$ give
\[
 \sum_i z_i h_i^2
   =\int_0^1\sum_{i\in I_t}h_i^2\,\dd t
   \le C\{z\log(en/z)+R\}.
\]
Thus, first by weighted Cauchy--Schwarz and then by Cauchy--Schwarz in the
observation law and Jensen's inequality,
\begin{align*}
 \E_{\mathbb Q}\sum_i z_i|h_i|
 &\le C\E_{\mathbb Q}\sqrt z\sqrt{z\log(en/z)+R}\\
 &\le C\sqrt D\sqrt{D\log(en/D)+\delta n+1+H}.
\end{align*}
Here $u\log(en/u)$ is increasing on $[0,n]$, allowing
$\E z\le D$ to be substituted.  Insert this in
\Cref{eq:stat-H-z}.  The term $C\sqrt{DH}$ is at most $H/2+C'D$;
the term $C\sqrt{D(\delta n+1)}$ is at most
$C'D+C'(\delta n+1)$.  Absorbing $H/2$ proves
\Cref{eq:stat-entropy}.  All estimates and choices were uniform over
$\nu$ and the cavities.
\end{proof}

\subsection{From entropy to a fixed-time density estimate}

The entropy bound gives a defective logarithmic Sobolev inequality.
Its energy coefficient grows slowly as the additive defect decreases.
Balancing these two terms in the hypercontractive estimate will give
the following bound.

\begin{proposition}[Stretched-exponential specific density]\label{prop:stat-density}
There are constants $C,c,d_*>0$ such that, for each fixed $d\ge d_*$,
on events of disorder probability tending to one,
\begin{equation}\label{eq:stat-density}
 \sup_{x,y}\log\frac{S_d(x,y)}{\mu(y)}
       \le Cn e^{-c d^{2/3}}.
\end{equation}
The assertion holds on one common event for the full system and all
single-site cavities.  The order of the quantifiers is fixed $d$, followed
by $n\to\infty$.
\end{proposition}

\begin{proof}
\emph{A defective logarithmic Sobolev inequality.}
We keep the dependence on the defect parameter explicit.  For
$0<\eta<1/2$, the concave function $u\sqrt{\log(en/u)}$ has at $u=\eta n$
a tangent with slope at most $\sqrt{\log(e/\eta)}$ and intercept at most
$\eta n/2$.  Thus \Cref{thm:stat-entropy}, with $\delta=\eta$, implies,
after normalizing $f^2$ to a density,
\begin{equation}\label{eq:stat-defective-lsi}
 \Ent_\mu(f^2)\le A_\eta\mathcal E_\mu(f)
                         +B_\eta\mu(f^2),
 \quad
 A_\eta=C(1+\sqrt{\log(1/\eta)}),\quad
 B_\eta=C(1+\eta n).
\end{equation}
The energy of the normalized density square root lies in $[0,n]$, as
required by \Cref{thm:stat-entropy}.  This also proves
\Cref{eq:stat-defective-lsi} for signed $f$, using
$\mathcal E_\mu(|f|)\le\mathcal E_\mu(f)$.

{\emergencystretch=3em
\emph{Hypercontractivity with the defect retained.}
Use the logarithmic-Sobolev-to-hyper\-contractivity argument of
Gross~\cite{Gross1975}, carrying the additive defect through the norm
calculation.  Let
$g_t=S_tf>0$, $p(0)=2$, and
$p'=4(p-1)/A_\eta$.  Reversibility gives\par}
\[
 \frac{\dd}{\dd t}\log\norm{g_t}_{p}
 =-\frac{\mathcal E_\mu(g_t^{p-1},g_t)}{\mu(g_t^p)}
    +\frac{p'}{p^2}\frac{\Ent_\mu(g_t^p)}{\mu(g_t^p)}.
\]
For $z,u>0$, Cauchy--Schwarz on the interval between them gives
\[
 (z^{p/2}-u^{p/2})^2
 \le\frac{p^2}{4(p-1)}(z^{p-1}-u^{p-1})(z-u).
\]
Apply \Cref{eq:stat-defective-lsi} to $g_t^{p/2}$ and use this inequality
edge by edge in the Dirichlet form.  The energy terms cancel, leaving
\[
 \frac{\dd}{\dd t}\log\norm{g_t}_{p(t)}
       \le B_\eta p'(t)/p(t)^2.
\]
Since $p(t)=1+\exp(4t/A_\eta)$ and
$\int_0^t p'/p^2\le1/2$, positivity and approximation give
\begin{equation}\label{eq:stat-hypercontractive}
 \norm{S_t}_{2\to p(t)}\le e^{B_\eta/2}.
\end{equation}

\emph{A pointwise kernel bound and the choice of defect.}
On $\norm J\le K$, every spin Hamiltonian has absolute value at most
$Kn/2$, so $\mu_{\min}\ge e^{-(K+\log2)n}$.  Consequently
\[
 \norm{S_t}_{2\to\infty}
 \le\exp\left(\frac{B_\eta}{2}+\frac{C n}{p(t)}\right).
\]
Self-adjointness and the semigroup property give
\begin{equation}\label{eq:stat-kernel-bound}
 \sup_{x,y}\frac{S_{2t}(x,y)}{\mu(y)}
 =\norm{S_{2t}}_{1\to\infty}
 \le\norm{S_t}_{2\to\infty}^2
 \le\exp\left(B_\eta+\frac{2Cn}{p(t)}\right).
\end{equation}
The equality is the elementary formula for the norm of a positive kernel
with respect to $\mu$.

We now choose the defect for the prescribed fixed time $d$.
Put $t=d/2$ and choose $\eta=\exp(-a d^{2/3})$, where $a>0$ is a
sufficiently small fixed constant.  For $d$ large,
$A_\eta\le C\sqrt a\,d^{1/3}$ after adjusting the lower threshold for
$d$.  Hence $p(d/2)^{-1}\le\exp(-c_1a^{-1/2}d^{2/3})$.
Choose $a$ so that $c_1a^{-1/2}>a$.  The logarithm in
\Cref{eq:stat-kernel-bound} is then at most
$C+Cn e^{-a d^{2/3}}$.  For this fixed $d$ its constant term is absorbed
once $n$ is sufficiently large, proving \Cref{eq:stat-density}.
All ingredients were common to the cavities, with their dimension $n-1$;
enlarging $C$ gives the stated common bound.
\end{proof}

\begin{remark}\label{rem:stat-order}
\Cref{prop:stat-density} gives an arbitrarily small positive specific
log-density cost by choosing a sufficiently large fixed burn-in.  It does
not assert a quenched estimate uniform in $d$ of order $\log n$.  The
later block argument uses precisely this dimension-first formulation.
\end{remark}
